% Round 16, global_height_audit. Independently reviewed by model_obstruction.
\section{Chow-correspondence adjunction and the abelian-motive height formula}
\label{r16ha:section}

The full mathematical source, including the supplied-alteration
and abelian-motive formulas, passed independent model-branch
review at SHA256
\path{09e3a701dcfb85898bd6a0c31a66d7ebe689fd1c18ad1bd502f7094dee76ba46}.
Only this activation notice differs from that reviewed source.

This section concerns the proved obstruction-zero domains
$\mathcal H$ and their canonical normalized real height. It makes
the correspondence compatibility explicit using the proper-intersection
height projection theorem, rather than assigning a pushforward to an
arbitrary actual form along a correspondence. The needed position
conditions are obtained over the ground number field.

\begin{lemma}[Moving with two maps, including all fiber dimensions]
\label{r16ha:moving}
Let $T,V,X$ be smooth projective geometrically integral varieties
over a number field $F$, and let $a:T\to V$, $b:T\to X$ be
morphisms. Put $\ell=\dim T$, and let $r,q\geq1$ satisfy
$r+q=\ell+1$, $r\leq\dim V$, $q\leq\dim X$.
Every pair of rational Chow classes $z\in\operatorname{CH}^r(V)$,
$w\in\operatorname{CH}^q(X)$ has representatives $Z,W$ such that
\[
 \begin{split}
 \dim a^{-1}|Z|&\leq\ell-r,\qquad
 \dim b^{-1}|W|\leq\ell-q,\\
 a^{-1}|Z|\cap b^{-1}|W|&=\varnothing.
 \end{split}
 \tag{HA16.1}\label{r16ha:position}
\]
The same assertion holds simultaneously for a finite list of pairs
of morphisms with common targets and the same complementary
source dimension. No extension of $F$ is needed for this movement.
\end{lemma}

\begin{proof}
For a proper morphism $a$, the sets on which its fiber dimension
equals $k$ are locally closed, by upper semicontinuity of fiber
dimension (Stacks 0D4I). Only finitely many $k$ occur. Decompose
each such set into finitely many irreducible locally closed strata
$S$, so that every fiber over $S$ has dimension $k(S)$.
The dimension formula gives
\[
                        \dim S+k(S)\leq\ell.
 \tag{HA16.2}\label{r16ha:dimension-bound}
\]
Indeed an irreducible component of the inverse image over the
generic point of $S$ attaining the fiber dimension has total
dimension $\dim S+k(S)$ and lies in $T$.

Apply the degree-zero part of Gorchinskiy's good moving-complex
Lemma~2.1 to the finite family of closures $\overline S\subset V$.
Its quasi-isomorphism is surjective on $H_0$, so it supplies a
cycle representative $Z$ of the given Chow class, with each
prime of $|Z|$ meeting every $\overline S$ properly. In particular,
\[
 \dim(|Z|\cap S)\leq\dim S-r,
 \qquad
 \dim a^{-1}(|Z|\cap S)
       \leq\dim S-r+k(S)\leq\ell-r.
\]
A negative upper bound means the intersection is empty.
The finite union over the strata proves the first bound in
\eqref{r16ha:position}. This argument also applies to strata of
the image when $a$ is not dominant; empty fibers contribute nothing.

Since $b$ is proper, the image
$B=b(a^{-1}|Z|)$ is a closed algebraic subset of $X$ and
$\dim B\leq\ell-r=q-1$. Apply the same good moving lemma on
$X$ to the closures of the fiber-dimension strata for $b$,
together with every irreducible component of $B$. It gives a
representative $W$ of $w$ meeting all those closed subvarieties
properly. The strata argument proves the second bound in
\eqref{r16ha:position}. For every component $B_j$ of $B$,
proper intersection with a codimension-$q$ cycle has dimension
at most $\dim B_j-q<0$, so $|W|\cap B=\varnothing$.
This is exactly the required inverse-support disjointness.

For a finite list of maps, take the finite union of the first
families when moving $Z$, and then the union of their closed
images $B$ and the second families when moving $W$. The same
dimension inequalities apply. Rational coefficients are handled
by clearing denominators and dividing after the movement.
\end{proof}

\begin{theorem}[Adjunction for every rational Chow correspondence]
\label{r16ha:adjunction}
Let $V/F$ and $X/F$ be smooth projective geometrically integral,
of dimensions $m,n$. Let
$\Gamma\in\operatorname{CH}^{m+d}(V\times_FX)_{\mathbf Q}$
be a rational Chow correspondence of degree $d$. Suppose
$z\in\mathcal H^r(V)$ and $w\in\mathcal H^q(X)$, with
\[
 1\leq r\leq m,\qquad 1\leq q\leq n,\qquad
                         r+d+q=n+1.
 \tag{HA16.3}\label{r16ha:degrees}
\]
Then the correspondence images belong to the required height
domains and
\[
 h_X^{\rm an}(\Gamma_*z,w)
              =h_V^{\rm an}(z,{}^t\Gamma_*w).
 \tag{HA16.4}\label{r16ha:adjunction-eq}
\]
This is an identity of the canonical global heights on Chow
classes. It includes singular correspondence cycles and maps
with arbitrary jumps in fiber dimension.
\end{theorem}

\begin{proof}
Membership of the two image classes follows from
Corollary~\ref{r16h:domain-functoriality}. Choose a cycle representing
$\Gamma$. By rational linearity it is enough to treat a prime
component. First suppose this component is geometrically integral.
Take a smooth projective resolution
$j:T\to\Gamma\subset V\times X$, with projections
$a:T\to V$, $b:T\to X$. Such a projective resolution is obtained
by the characteristic-zero embedded resolution theorem on the
smooth projective ambient $V\times X$; see W\l odarczyk,
\emph{Simple Hironaka resolution in characteristic zero},
Theorem~1.0.2. Its restriction to the strict transform is
projective and generically of degree one. Here
$\ell=\dim T=n-d$, so \eqref{r16ha:degrees} gives
$r+q=\ell+1$ and in particular $r,q\leq\ell$.

Choose representatives satisfying Lemma~\ref{r16ha:moving} and put
\[
                 A=a^!Z,\qquad B=b^!W.
 \tag{HA16.5}\label{r16ha:refined-cycles}
\]
The expected inverse-image dimension bounds make these the
specified top-dimensional refined cycles on their respective
inverse supports. Their coefficients are the refined Gysin
coefficients. They are not replaced by those of a reduced
inverse image. In Chow groups they represent $a^*z,b^*w$,
which belong to $\mathcal H^r(T),\mathcal H^q(T)$ by the
proved pullback property. Their supporting envelopes are
disjoint by \eqref{r16ha:position}.

Apply Theorem~\ref{r16hp:height-projection} to $b:T\to X$,
the source cycle $A$, and the target cycle $W$. Its numerical
condition is $r+q=\ell+1$; its proper-intersection condition
is the second bound in \eqref{r16ha:position}; and its disjointness
condition follows from the last condition there. It gives
\[
                  h_X^{\rm an}(b_*A,W)=h_T^{\rm an}(A,B).
 \tag{HA16.6}\label{r16ha:first-projection}
\]
Next apply that theorem to $a:T\to V$, the source cycle $B$,
and target cycle $Z$. Now use the first dimension bound and
the same disjointness. With the already proved symmetry, this gives
\[
                  h_T^{\rm an}(A,B)=h_V^{\rm an}(Z,a_*B).
 \tag{HA16.7}\label{r16ha:second-projection}
\]
These applications use the proved all-actual proper-slice
residue for the concrete obstruction-zero choices and the
precise classical logarithmic Green pullback at infinity.
They do not pull back an arbitrary complex current or infer
an excess-intersection formula outside the displayed position.

Since $j_*[T]=[\Gamma]$, algebraic refined projection gives
$[b_*A]=\Gamma_*z$ and $[a_*B]={}^t\Gamma_*w$.
Equations \eqref{r16ha:first-projection} and
\eqref{r16ha:second-projection}, together with Chow independence
of the heights, prove \eqref{r16ha:adjunction-eq}.
For finitely many geometrically integral components, either use
the simultaneous movement in the lemma or sum the individual
Chow identities.

For a prime correspondence which is not geometrically integral,
choose a finite extension $E/F$ over which its finitely many
geometric components are defined and geometrically integral.
Characteristic zero makes its flat base change generically
reduced. Apply the preceding argument to every component with
its flat cycle coefficient. Chow pullback, proper push and
correspondence composition commute with this flat base change.
The source classes remain in $\mathcal H$ after extension by
the native constant-extension theorem, and the normalized
heights on their pullbacks are unchanged. Thus the identity
over $E$ is exactly the identity over $F$; no unrecorded
$[E:F]$ factor is introduced. Finally linearity handles every
rational representative of the correspondence class.
\end{proof}

\begin{corollary}[Scalar comparison through a supplied admissible alteration]
\label{r16ha:alteration-comparison}
Let $X/F$ be smooth projective geometrically integral of dimension
$n$, and let $z,w$ be geometrically homologically trivial rational
Chow classes of codimensions $r,s\geq1$, with $r+s=n+1$.
Suppose that for some finite extension $E/F$ there is a
generically finite dominant morphism
\[
                      f:Y\longrightarrow X_E
\]
of degree $d$, with $Y/E$ smooth projective geometrically integral,
and a supplied regular flat projective model
$\mathcal Y/\mathcal O_E$ on which $f^*z_E$ and $f^*w_E$
admit numerically admissible rational Chow lifts $\Theta,\Phi$.
Then $z\in\mathcal H^r(X)$, $w\in\mathcal H^s(X)$, and
\[
 h_X^{\rm an}(z,w)
       =\frac1d h_Y^{\rm an}(f^*z_E,f^*w_E)
       =\frac1{d[E:\mathbf Q]}
                     \widehat{\deg}(\widehat\Theta\cdot\widehat\Phi).
 \tag{HA16.10}\label{r16ha:alteration-height}
\]
The arithmetic lifts on the right have zero archimedean curvature.
The right side is independent of the supplied alteration, model
and admissible lifts whenever those data satisfy these hypotheses.
No extension of $f$ to a morphism of the chosen arithmetic models
is required.
\end{corollary}

\begin{proof}
The numerical model theorem gives
$f^*z_E\in\mathcal H^r(Y)$ and
$f^*w_E\in\mathcal H^s(Y)$. At every finite place of $E$,
the actual obstruction of a pulled class is the pullback of the
target obstruction. The all-degree dominant Bott--Chern
injectivity \texttt{r16pb:dominant} makes the latter zero.
Thus $z_E,w_E$ lie in the two height domains on $X_E$.
Finite constant-extension injectivity, with the native actual
sheaf and complete field degrees, then gives the stated
memberships over $F$. This is a descent of obstruction zero;
it is not a converse inferred from a bare strong Green form.

Apply Theorem~\ref{r16ha:adjunction} to the graph correspondence
of $f$, with source class $f^*z_E$ and target class $w_E$.
The algebraic projection formula gives
$f_*f^*z_E=d z_E$, since $f_*[Y]=d[X_E]$.
Consequently
\[
                 h_Y^{\rm an}(f^*z_E,f^*w_E)
                         =d\,h_{X_E}^{\rm an}(z_E,w_E).
\]
Native number-field invariance identifies the last height with
$h_X^{\rm an}(z,w)$. Theorem~\ref{r16mm:global}, applied only
on the supplied model of $Y$, gives the second equality in
\eqref{r16ha:alteration-height}. Every set of supplied data
produces the same canonical left side, which proves the
independence assertion.

This argument does not use an arithmetic-model projection
formula for $f$. Its degree calculation is the ordinary generic
Chow identity followed by the already proved analytic height
adjunction. In particular it never assumes an alteration is
finite outside the specified cycle. Existence of such an
admissible model lift for general classes remains open.
\end{proof}

\begin{corollary}[An explicit height formula for a supplied abelian Chow factorization]
\label{r16ha:motivic-formula}
Let $X/F$ be smooth projective geometrically integral of
dimension $n$. Suppose that after a finite extension $E/F$
there is a rational Chow identity
\[
                 [\Delta_{X_E}]=\sum_j v_j\circ u_j,
 \tag{HA16.8}\label{r16ha:factorization}
\]
where $u_j:X_E\rightsquigarrow A_j$ has degree $d_j$,
$v_j:A_j\rightsquigarrow X_E$ has degree $-d_j$, and
$A_j/E$ is an abelian variety of dimension $a_j$.
For every pair of geometric homological classes
$z\in\operatorname{CH}^r(X)_{\mathbf Q}$,
$w\in\operatorname{CH}^{n+1-r}(X)_{\mathbf Q}$, one has
\[
 h_X^{\rm an}(z,w)
   =\sum_j h_{A_j}^{\rm an}
                   (u_{j*}z_E,{}^tv_{j*}w_E)
   =\frac1{[E:\mathbf Q]}\sum_j
        \bigl\langle u_{j*}z_E,{}^tv_{j*}w_E\bigr\rangle_{A_j,\mathrm{Kun}}.
 \tag{HA16.9}\label{r16ha:motivic-height}
\]
The two codimensions in the $j$-th term are
$r+d_j$ and $a_j+1-r-d_j$. Terms outside the Chow range, or
in codimension zero, are zero on the geometric homological
kernel and are omitted. The last pairing is K\"unnemann's
published classical pairing in the fixed scalar sign convention.

The value is independent of the supplied factorization and its
finite field of definition. No regular arithmetic model of
$X$ is required for this formula.
\end{corollary}

\begin{proof}
Theorem~\texttt{r16ab:motivic} puts both complete geometric
homological groups in $\mathcal H$, including after the indicated
finite extension. Native field invariance lets us work over $E$.
Apply \eqref{r16ha:factorization} to $z_E$, use rational
bilinearity, and apply Theorem~\ref{r16ha:adjunction} to each
$v_j$. This gives the first equality in
\eqref{r16ha:motivic-height}. Geometric homological triviality
is preserved by all the correspondences, so their abelian
image classes lie in the full abelian height domains.
The second equality is Corollary~\ref{r16mm:abelian}.
If a degree leaves the Chow range the corresponding image
vanishes; if it is zero its geometric cycle class detects its
rank, so its geometric homological image is zero as well.

Every factorization produces the same previously defined
canonical analytic height on the left. This proves independence
of the factorization and field without imposing a uniqueness
axiom on arbitrary scalar pairings. The comparison invokes
regular admissible models only for the abelian varieties,
where the published theorem supplies them after finite constant
extension. It is not a general arithmetic alteration comparison
for $X$ and is not a new classical abelian existence theorem.
\end{proof}

\paragraph{Primary and internal inputs.}
The finite-family movement is Gorchinskiy, Lemma~2.1,
\url{https://arxiv.org/pdf/0802.1437}. Fiber-dimension
upper semicontinuity is
\url{https://stacks.math.columbia.edu/tag/0D4I}.
Projective resolution is obtained from the embedded theorem
in W\l odarczyk's published paper, Theorem~1.0.2,
\url{https://www.math.purdue.edu/~wlodarcz/singularities/Resolution.pdf}.
The analytic inputs are the independently reviewed
\path{global_height/proper_intersection_extension.tex},
\path{global_height/obstruction_zero_reciprocity_and_height.tex},
and the separate abelian comparison in
\path{global_height/regular_admissible_model_comparison.tex}.
